% Fragment partage : inclus par compte-certifie-officiel.tex (dossier autonome)
% et par dossier-conception-alanya.tex (dossier client assemble).
% Ne pas y mettre de preambule ni de \part : c'est l'appelant qui les pose.

\chapter{Ce que la certification dit, et ne dit pas}

\begin{alerte}[Révisé par le volet 8]
L'offre payante est désormais arrêtée. Hors phase gratuite, la coche personnelle demande un
abonnement actif \emph{en plus} de la vérification, et elle devient accessible à tout abonné
--- plus seulement aux personnalités publiques. Le dossier, le coffre et la purge décrits
ici restent valables~; les tables d'abonnement de ce volet sont remplacées par celles du
volet~8.
\end{alerte}

\section{Une seule question}

Le badge répond à \emph{une} question~: \guillemotleft~ce compte est-il bien qui il prétend
être~?~\guillemotright{} Rien d'autre.

Il ne dit pas que le commerce est sérieux, que ses prix sont justes, ni qu'il livre à
l'heure. C'est une attestation d'identité, pas une recommandation --- et la communication
doit le refléter, sous peine d'engager Alanya sur ce qu'elle ne vérifie pas.

\begin{center}
\begin{tabularx}{\linewidth}{@{}p{4.6cm}X@{}}
\toprule
\thd{Ce qu'on vérifie} & Que l'entreprise existe, et que la personne qui tient le compte est
habilitée à la représenter. \\
\thd{Ce qu'on ne vérifie pas} & La qualité des produits, l'honnêteté commerciale, le respect
des délais. \\
\bottomrule
\end{tabularx}
\end{center}

\section{Qui peut être vérifié}

Les deux axes du socle étant indépendants, la vérification ne concerne pas que les
commerces~:

\begin{center}
\begin{tabular}{@{}llp{6.4cm}@{}}
\toprule
\thd{Genre} & \thd{Rendu} & \thd{Cas typique} \\
\midrule
Business  & \badgePastillePanier[1] & Une entreprise, une boutique, un cabinet \\
Personnel & \badgePastilleCoche[1]  & Une personnalité publique, un journaliste \\
\bottomrule
\end{tabular}
\end{center}

\chapter{Le parcours du demandeur}

\section{Trois états}

\begin{center}
  \imageOuCadre{maquettes/certification-parcours.png}{\linewidth}%
    {Dépôt du dossier, attente, compte vérifié avec relance à J-30\\
     \texttt{maquettes/certification-parcours.png}}
\end{center}

\begin{note}[Maquette interactive]
\href{https://claude.ai/code/artifact/fb9bb169-1370-4cad-be00-0101a68ba2fe}%
{\texttt{claude.ai/code/artifact/fb9bb169}}\\
Source versionnée~: \code{docs/architecture/maquettes/certification.html}
\end{note}

\section{Les pièces demandées}

Proposition, à valider --- elle est calquée sur le contexte camerounais et ouest-africain~:

\begin{center}
\begin{tabularx}{\linewidth}{@{}p{4.8cm}p{2.2cm}X@{}}
\toprule
\thd{Pièce} & \thd{Statut} & \thd{Pourquoi} \\
\midrule
Registre de commerce (RCCM) & obligatoire & Prouve que l'entreprise existe \\
Pièce d'identité du gérant & obligatoire & Relie la personne à l'entreprise \\
Justificatif d'adresse & facultatif & Corrobore le lieu déclaré \\
\bottomrule
\end{tabularx}
\end{center}

Pour un compte \emph{personnel} vérifié, la liste change~: pièce d'identité, plus un
élément de notoriété publique --- page officielle, article de presse, site institutionnel.

\begin{note}[Annoncer la rétention dès le dépôt]
L'écran de dépôt dit, avant l'envoi, que les pièces sont chiffrées, supprimées 90~jours
après la décision, et que chaque consultation est journalisée. Le dire après serait perçu
comme une justification~; le dire avant est ce qui fait déposer.
\end{note}

\section{Ce qui est promis, et ce qui ne l'est pas}

\guillemotleft~Réponse sous 5~jours ouvrés~\guillemotright{} plutôt qu'un compte à rebours.
Un délai affiché à la minute transforme le sixième jour en manquement~; une fourchette
laisse la marge que l'instruction demande réellement.

\chapter{Le parcours de l'administrateur}

\section{L'écran d'instruction}

\begin{center}
  \imageOuCadre{maquettes/certification-instruction.png}{\linewidth}%
    {Instruction d'un dossier : identité déclarée, pièces, échéance de purge\\
     \texttt{maquettes/certification-instruction.png}}
\end{center}

À gauche ce que le compte déclare, à droite ce qu'il apporte comme preuve. L'administrateur
compare les deux~; c'est tout le travail.

\section{Trois issues, pas deux}

\begin{center}
\begin{tabularx}{\linewidth}{@{}p{4.4cm}X@{}}
\toprule
\thd{Approuver} & Le badge apparaît, l'échéance est posée. \\
\thd{Refuser avec motif} & Le motif est \textbf{obligatoire} et repart dans la notification.
Sans lui, le demandeur redépose le même dossier incomplet et la file se remplit de
doublons. \\
\thd{Demander une pièce} & Le dossier \emph{reste} en attente, le demandeur est prévenu de
ce qui manque. \\
\bottomrule
\end{tabularx}
\end{center}

\begin{note}[La troisième issue manquait]
Ni \code{alanya-contexte-conception.md} ni \code{ideas/todo.tex} ne la mentionnaient. C'est
pourtant le cas le plus fréquent en pratique~: le dossier n'est pas mauvais, il est
incomplet. Sans cette issue, l'administrateur doit refuser puis expliquer --- ce qui est
vécu comme un rejet.
\end{note}

\section{Qui décide}

\textbf{Tout administrateur} (\code{type\_compte >= 1}). C'est cohérent avec l'existant~:
bannir et débannir sont déjà ouverts à tout administrateur
(\code{src/routes/admin.js}). Chaque décision garde son auteur dans \code{reviewed\_by} et
sa date --- la traçabilité remplace la restriction.
